\subsection{No repair by an invisible direct summand}

An added generically acyclic summand cannot repair a
global graph part carrying the required generic
cohomology. This conclusion assumes an actual derived
direct-sum decomposition; it does not split arbitrary
extensions.

Let \(A\) be a positively graded polynomial algebra
over a characteristic-zero field \(k\), with \(A_0=k\).
Let \(\tau\) be a nontrivial graded involution.
Write \(A_e,A_\tau\) for the identity and twisted
graph bimodules, and \(\mathfrak m=A_{>0}\) for the
homogeneous origin, preserved by \(\tau\).
Work with bounded graded bimodule complexes perfect
on both sides and with finite graded cohomology on
both sides. All scalar complexes below are bounded
complexes of finite graded \(k\)-spaces.
Put \(Q=\operatorname{Frac}(A)\).

\begin{finding}[\fid{F59}, the invisible summand boundary]
\label{f:invisibleSummandNoGo}
Suppose there is a graded derived decomposition
\begin{equation}\label{eq:invisibleSummandDecomposition}
 X=(U\otimes_kA_e)\oplus(V\otimes_kA_\tau)\oplus Z,
\end{equation}
where \(U,V\) have nonzero cohomology and
\[
                         Q\otimes_A^{\mathbf L}Z=0.
\]
If a scalar complex \(C\) gives an actual isomorphism
\begin{equation}\label{eq:invisibleSummandQuadratic}
 X\otimes_A^{\mathbf L}X\simeq C\otimes_kX,
\end{equation}
then \(Z=0\) in the derived category.

For the actual ambient generator and scalar Euler
class in Finding~\ref{f:alteredScalarSoergel}, no
decomposition of this kind can therefore repair the
quadratic: after \(Z=0\), its cohomology is graded
left-free and that finding applies.
\end{finding}

\begin{proof}
\par\smallskip\noindent\textbf{Generic dimensions.}\quad
Forget internal grading and set
\[
 u(t)=\sum_i\dim_kH^i(U)t^i,\qquad
 v(t)=\sum_i\dim_kH^i(V)t^i,\qquad
 d(t)=\sum_i\dim_kH^i(C)t^i.
\]
After left localization, \(X\) is a sum of the two
graph complexes and has its canonical right \(Q\)-action.
Its tensor with \(X\) is therefore its tensor over
\(Q\) with \(Q\otimes_A^{\mathbf L}X\).
This justifies localization of the square; no
monoidality of arbitrary double-localized bimodules
is assumed.

Since \(\tau^2=e\), generic K\"unneth gives
\[
                  u^2+v^2=du,\qquad 2uv=dv.
\]
The polynomial \(v\) is nonzero, hence \(d=2u\).
Then \(v^2=u^2\), and the Laurent polynomial domain,
with \(u+v\ne0\), gives \(v=u\).
Thus, for the positive integer \(h=u(1)=v(1)\),
\begin{equation}\label{eq:invisibleScalarDimensions}
 \dim H^*(U)=\dim H^*(V)=h,\qquad
                         \dim H^*(C)=2h.
\end{equation}
Only these ungraded dimensions are used.
No equality of full internal shift data is inferred.

\par\smallskip\noindent\textbf{A positive derived-fiber size.}\quad
For a displayed bimodule complex \(M\), define
\begin{equation}\label{eq:derivedOriginSize}
 \mu(M)=\sum_i\dim_k
             H^i(k\otimes_A^{\mathbf L}M),
 \qquad k=A/\mathfrak m.
\end{equation}
The fiber uses the left action and then forgets the
remaining right action.
Left perfection makes this a bounded complex of
finite-dimensional vector spaces. Hence \(\mu\) is
a finite nonnegative integer, additive on direct sums
and invariant under derived isomorphisms.

It detects nonzero objects in this graded category.
If \(r\) is the highest degree with \(H^r(M)\ne0\),
top-degree truncation gives
\[
 H^r(k\otimes_A^{\mathbf L}M)
       =k\otimes_AH^r(M)\ne0.
\]
Higher Tor contributes only in lower degrees, and the
last nonvanishing is graded Nakayama for a nonzero
finite graded module. Thus \(\mu(M)=0\) implies \(M=0\).

For a scalar complex \(W\), K\"unneth over \(k\) gives
\begin{equation}\label{eq:derivedOriginScalarSize}
 \mu(W\otimes_kM)=\dim_kH^*(W)\,\mu(M).
\end{equation}
The two graph objects have size one.
Right tensoring with \(A_\tau\) leaves the underlying
left module unchanged up to its graph identification.
Left tensoring twists the left action by \(\tau\)
or \(\tau^{-1}\), which preserves \(\mathfrak m\).
Consequently
\begin{equation}\label{eq:derivedOriginGraphSize}
 \mu(A_\tau\otimes_A^{\mathbf L}M)
   =\mu(M\otimes_A^{\mathbf L}A_\tau)=\mu(M).
\end{equation}

All the tensor products used here remain left-perfect,
so their sizes are finite. Indeed, build the second
factor as a left-perfect complex from finitely many
shifted copies of \(A\) using cones and summands.
Tensor with the first factor to obtain the same finite
construction from that left-perfect factor.
The corresponding right-perfect argument keeps the
products in the stated bimodule category.

\par\smallskip\noindent\textbf{The dimension contradiction.}\quad
Set \(Y=(U\otimes A_e)\oplus(V\otimes A_\tau)\).
The four graph products and
\eqref{eq:invisibleScalarDimensions} give
\[
 \mu(Y\otimes_A^{\mathbf L}Y)=4h^2
                             =\mu(C\otimes_kY).
\]
The cross terms have sizes
\[
 \mu(Y\otimes_A^{\mathbf L}Z)
    =\mu(Z\otimes_A^{\mathbf L}Y)
    =2h\mu(Z),\qquad
 \mu(C\otimes_kZ)=2h\mu(Z).
\]
Apply \(\mu\) to \eqref{eq:invisibleSummandQuadratic}
and expand the actual direct sums.
Canceling the common numerical term \(4h^2\) leaves
\[
                     2h\mu(Z)
                       +\mu(Z\otimes_A^{\mathbf L}Z)=0.
\]
Both terms are nonnegative and \(h>0\).
Hence \(\mu(Z)=0\), and the proved derived graded
Nakayama implication gives \(Z=0\).
\end{proof}

No cancellation theorem for objects, Krull--Schmidt
premise or bimodule formality is used. The invariant
is the positive total cohomology dimension of a
derived fiber, not its Euler characteristic.
The theorem excludes an invisible correction supplied
as a direct summand. It does not assert that a general
coupled complex decomposes into a free graph part and
such a correction. Nonsplit extensions with nonflat
cohomology remain outside this argument.

% Separate paragraph for Part III.
\paragraph{The invisible summand test in Grothendieck's language.}
The affine module/sheaf framework is EGA~I,
Theorem~1.3.7, p.~86, Corollary~1.3.8, p.~87,
and Theorem~1.4.1, pp.~90--92 \cite{EGA}.
Finite perfect complexes and derived fibers use
SGA~6, I, Definition~4.7, pp.~121--122, and
Example~4.8, p.~122 \cite{SGA6}.
The grading has the split-torus action framework
of SGA~3, I, 4.7.2--3 and cote~90, scan~17;
the effective-dimension versus virtual-Euler distinction
has the symmetric coefficient language of SGA~6,
V, \S\S2--3 and cote~133, scan~18
\cite{SGA3,SGA6,Archive}.

Every derived-fiber finiteness and nonvanishing
implication is proved with its actual hypotheses.
No general cohomology-and-base-change claim replaces
it. The dictionary introduces no \(\lambda\)-ring on
the noncommutative word algebra, source canonical
lattice, arithmetic Frobenius or Weil-purity claim.
